% Vietnamese translation for OI Public Library
% Source-PDF: IOI中国国家候选队论文/2002/何林.pdf
\documentclass[11pt]{article}
\usepackage{tikz}
\usepackage{oipl-vn}

\title{Phỏng đoán và ứng dụng của nó}
\author{He Lin}
\date{}

\begin{document}
\maketitle

\origtitle{Conjecture and its applications}
\sourcepath{IOI China candidate team papers / 2002 / He Lin.pdf}

\section*{Tổng quan}

Phỏng đoán là một chiến lược giải bài quan trọng, rất có ích cho việc rèn
luyện năng lực tư duy tổng hợp.

Phần đầu của bài viết giới thiệu các yêu cầu cơ bản và các bước thực hiện
phỏng đoán, đồng thời nêu hai phương pháp quan trọng để đi tới phỏng đoán:
tương tự và quy nạp.

Phần thứ hai và thứ ba lần lượt trình bày các tính chất cơ bản, định nghĩa và
các bước giải bài nói chung của tương tự và quy nạp. Mỗi phần đều dùng một số
bài toán cụ thể để minh họa ứng dụng của phỏng đoán trong thi tin học.

Cuối cùng, bài viết tổng kết toàn văn và trình bày một cách toàn diện về
phương pháp, định nghĩa, phân loại và tầm quan trọng của phỏng đoán.

\paragraph{Từ khóa.}
Phỏng đoán; tương tự; quy nạp.

\tableofcontents

\section{Dẫn nhập}

Phỏng đoán là phương pháp cơ bản và quan trọng nhất để giải quyết vấn đề, là
công cụ mạnh để thăm dò từ thế giới đã biết sang thế giới chưa biết. Trong
thi tin học, khi đề bài ngày càng che giấu sâu hơn bản chất của nó, phương
pháp nghiên cứu bằng ``phỏng đoán'' cũng giữ vai trò ngày càng quan trọng.

Phỏng đoán không phải là đoán bừa. Ta phải phân tích tổng hợp các điều kiện
để đưa ra một suy đoán ``có lý'', rồi dưới sự dẫn dắt của suy đoán ấy tiến
hành luận chứng ``có căn cứ''. Vì vậy phỏng đoán tuyệt đối không chỉ là may
mắn hay đầu cơ. Cũng như một học sinh tiểu học rất khó nêu ra định lý lớn
Fermat, nếu không quan sát và phân tích sâu thì ta cũng không thể đưa ra bất
kỳ phỏng đoán nào có giá trị.

Các bước cơ bản khi phỏng đoán là:
\[
\begin{array}{ccccc}
\text{Đề xuất phỏng đoán}
&\longrightarrow&
\text{kiểm tra bằng dữ liệu nhỏ}
&\xrightarrow{\text{đúng}}&
\text{chứng minh kết luận}
\\
&&\bigg\downarrow\scriptstyle{\text{sai}}&&\bigg\downarrow\\
&&\text{sửa phỏng đoán}
&\longleftarrow&
\text{giải quyết vấn đề}
\end{array}
\]

Do đó, phỏng đoán là một quá trình dựa trên phân tích sâu, liên tục thăm dò
và lặp lại việc sửa chữa. Nó không thể hoàn thành trong một bước. ``Đề xuất
phỏng đoán'' là phần lõi của toàn bộ quá trình. Các phương pháp chính để có
được phỏng đoán là tương tự và quy nạp.

\section{Phỏng đoán bằng tương tự}

Tương tự là một hình thức suy luận: từ một số tính chất giống hoặc gần giống
nhau của hai đối tượng, ta suy ra rằng chúng cũng có thể giống hoặc gần giống
nhau ở những tính chất khác.

Biểu diễn hình thức như sau: sự vật \(A\) có các tính chất \(a,b,c,d\); sự vật
\(B\) có các tính chất tương ứng \(a',b',c'\). Từ đó ta phỏng đoán rằng \(B\)
cũng có tính chất \(d'\) tương ứng với \(d\).

Tương tự là một suy luận hợp lý nhưng chủ quan và chưa đầy đủ. Vì vậy, muốn
xác nhận tính đúng đắn của phỏng đoán nhận được bằng tương tự, ta vẫn phải
qua một chứng minh logic nghiêm ngặt.

\subsection{Tương tự với bài toán quen thuộc}

Nếu bài toán \(A\), là bài toán quen thuộc, đã được giải, còn bài toán \(B\),
là bài toán lạ, giống \(A\) về hình thức, cấu trúc hoặc tính chất, thì có thể
phỏng đoán rằng cách giải \(B\) tương tự cách giải \(A\). Có thể hình dung:
\[
\begin{array}{ccc}
\text{bài toán }A & \xrightarrow{\text{tương tự}} & \text{bài toán }B\\
\downarrow\scriptstyle{\text{giải}} && \downarrow\scriptstyle{\text{giải}}\\
\text{thuật toán }A & \xrightarrow{\text{phỏng đoán}} & \text{thuật toán }B
\end{array}
\]

Các bước cơ bản khi tương tự với một bài toán quen thuộc là:
\begin{enumerate}
  \item \term{Phân tích.} Phân tích đặc trưng của bài toán và trừu tượng hóa
  thành mô hình.
  \item \term{Liên tưởng.} So sánh với bài toán quen thuộc có đặc trưng hoặc
  mô hình tương tự.
  \item \term{So sánh.} Sau khi xác định đối tượng tương tự, so sánh hai bài
  toán và phân tích điểm giống, điểm khác.
  \item \term{Phỏng đoán.} Dựa vào bài toán đã biết để phỏng đoán cách giải
  bài toán mới.
\end{enumerate}

\subsubsection*{Ví dụ dẫn nhập: người khổng lồ và ma}

\paragraph{Mô tả bài toán.}
Trên mặt phẳng có \(n\) người khổng lồ và \(n\) con ma, với
\(1\le n\le 50\). Mỗi người khổng lồ có vũ khí laser kiểu mới, có thể tấn
công ma theo một đường thẳng. Tuy nhiên vũ khí này có một nhược điểm lớn:
laser bắn ra từ hai vũ khí khác nhau không được cắt nhau, nếu không sẽ có
nguy cơ phát nổ.

Biết rằng không có ba cá thể bất kỳ, kể cả người khổng lồ và ma, cùng nằm
trên một đường thẳng. Mỗi người khổng lồ phụ trách tấn công một con ma, và
dĩ nhiên các tia laser không được cắt nhau.

Yêu cầu là tìm một phương án tấn công bất kỳ, tức xác định người khổng lồ nào
tấn công con ma nào.

\paragraph{Dữ liệu vào.}
Dòng đầu chỉ có một số nguyên \(n\). Tiếp theo là \(n\) dòng, mỗi dòng có hai
số nguyên \(x,y\), biểu diễn tọa độ người khổng lồ. Sau đó là \(n\) dòng,
mỗi dòng có hai số nguyên \(x,y\), biểu diễn tọa độ ma.

\paragraph{Dữ liệu ra.}
Tệp ra có \(n\) dòng. Dòng thứ \(i\) chứa một số nguyên \(p\), nghĩa là người
khổng lồ thứ \(i\) tấn công con ma thứ \(p\).

\paragraph{Phân tích.}
Thực chất bài toán yêu cầu xác định một song ánh giữa người khổng lồ và ma,
điều này rất tự nhiên gợi đến bài toán ghép cặp.

Trừu tượng hóa người khổng lồ và ma thành các đỉnh. Gọi tập đỉnh người khổng
lồ là \(X\), tập đỉnh ma là \(Y\). Nối một cạnh giữa mọi \(x\in X\) và
\(y\in Y\), rồi dùng thuật toán Hungary để tìm một ghép hoàn hảo của đồ thị
hai phía.

Nhưng bài toán còn có một điều kiện đặc biệt: hai cạnh ghép bất kỳ không được
cắt nhau. Vì thế ghép cặp đơn giản chưa đủ. Ta giả sử có hai cạnh ghép cắt
nhau và xét tính chất của chúng.

Nếu \(A\) tấn công \(D\), còn \(C\) tấn công \(B\), thì hai cạnh ghép \(AD\)
và \(CB\) cắt nhau tại \(O\). Một ý tưởng tự nhiên là đổi thành phương án
\(A\) ghép với \(B\), \(C\) ghép với \(D\). Quan sát hình học cho thấy
\[
  OA+OB>AB,\qquad OC+OD>CD,
\]
nên
\[
  AD+BC>AB+CD.
\]
So với phương án cũ, phương án mới làm ``tổng độ dài các cạnh ghép'' giảm đi.
Vì vậy, mỗi ghép hoàn hảo có chứa các cạnh cắt nhau đều có thể chuyển thành
một ghép hoàn hảo có tổng độ dài cạnh nhỏ hơn. Nói cách khác, ghép hoàn hảo
có tổng độ dài cạnh nhỏ nhất chắc chắn không chứa hai cạnh cắt nhau.

Ta chuyển mô hình ghép hoàn hảo ban đầu thành mô hình ghép tối ưu trên đồ thị
hai phía: trọng số mỗi cạnh là khoảng cách Euclid giữa hai đầu mút. Bài toán
được giải.

\paragraph{Nhận xét.}
Các đỉnh trong đồ thị hai phía của bài này có tính đặc biệt: mỗi đỉnh tương
ứng với một điểm trên mặt phẳng. Vì vậy giữa các đỉnh không chỉ có quan hệ
logic, mà còn có quan hệ trực quan hình học. Chính nhờ tận dụng đầy đủ quan hệ
ẩn ``trực quan hình học'' này mà mô hình thuật toán trở nên rõ ràng.

\subsubsection*{Ví dụ 1: nỗi phiền của ếch}

\paragraph{Mô tả bài toán.}
Trong ao có \(n\) lá sen, với \(1\le n\le 1000\), vừa đúng tạo thành một đa
giác lồi. Ta đánh số các lá sen theo chiều kim đồng hồ là
\(1,2,\ldots,n\).

Một chú ếch nhỏ đứng trên lá sen số 1. Nó muốn nhảy qua mỗi lá sen đúng một
lần, và có thể từ lá đang đứng nhảy tới bất kỳ lá nào khác. Đồng thời nó muốn
tổng quãng đường đã nhảy là ngắn nhất. Hãy lập trình giúp ếch thiết kế một
lộ trình.

\paragraph{Dữ liệu vào.}
Dòng đầu có một số nguyên \(n\). Tiếp theo là \(n\) dòng, mỗi dòng có hai số
nguyên \(x,y\), biểu diễn tọa độ lá sen thứ \(i\).

\paragraph{Dữ liệu ra.}
Tệp ra chỉ có một số nguyên, là tổng quãng đường ngắn nhất.

\paragraph{Phân tích.}
Đây là một bài khó.

Trước hết trừu tượng hóa mỗi lá sen thành một đỉnh; giữa hai đỉnh bất kỳ
\(x,y\) nối một cạnh, và trọng số cạnh là khoảng cách Euclid giữa \(x\) và
\(y\). Dễ thấy bản chất của bài toán là tìm một đường Hamilton xuất phát từ
1, vốn là một bài toán NP-khó kinh điển. Vì \(n\) có thể tới 1000, trước hết
ta loại bỏ cách tìm kiếm.

Ta tự nhiên nêu phỏng đoán đầu tiên:

\begin{quote}
\term{Phỏng đoán 1.} Bắt đầu từ 1, lần lượt đi qua mọi điểm theo chiều kim
đồng hồ sẽ cho đường đi ngắn nhất.
\end{quote}

Nhưng một phản ví dụ đơn giản đã bác bỏ nó. Lấy \(ABCD\) là một hình bình
hành. Theo phỏng đoán 1, đường tốt nhất phải là \(ABCD\); nhưng đường ngắn
nhất thật sự là \(ABDC\). Vì vậy việc đi đơn giản theo thứ tự trên bao lồi
không được.

Ta tiếp tục cân nhắc kỹ tính chất của bài toán. Đồ thị được xây dựng có ba
đặc điểm:
\begin{enumerate}
  \item là một đồ thị đầy đủ;
  \item trọng số cạnh bằng khoảng cách Euclid giữa hai đầu mút;
  \item mỗi đỉnh trong đồ thị tương ứng với một điểm trên mặt phẳng.
\end{enumerate}
Những tính chất này tự nhiên gợi tới ví dụ dẫn nhập. Đó là sự giống nhau bề
mặt.

Vì mỗi đỉnh trong đồ thị tương ứng với một điểm trên mặt phẳng, giữa các đỉnh
không chỉ có quan hệ logic do cạnh nối tạo ra, mà còn có quan hệ trực quan
hình học đặc biệt. Tận dụng điều kiện ẩn này là chìa khóa giải bài, và điều
đó trùng khớp với đặc trưng cơ bản của ví dụ dẫn nhập. Đó là sự giống nhau
bản chất.

Khi giải ví dụ dẫn nhập, ta dùng tính chất này để liên tục giảm ``tổng độ dài
các cạnh ghép''. Bài này cũng là một bài toán đường đi ngắn nhất. Dựa trên
hai điểm tương tự ấy, ta phỏng đoán:

\begin{quote}
\term{Phỏng đoán 2.} Trong đường Hamilton ngắn nhất không tồn tại hai cạnh
cắt nhau.
\end{quote}

\paragraph{Chứng minh.}
Giả sử từ 1, đỉnh đầu tiên được đi tới là \(A\), với \(2<A<n\). Vì mọi điểm
nằm trên một bao lồi và \(2<A<n\), nên ngoài hai đỉnh \(1,A\), các đỉnh còn
lại chắc chắn bị đoạn thẳng \(\langle 1,A\rangle\) chia thành hai phần ở hai
phía của đoạn ấy. Không mất tính tổng quát, giả sử từ \(A\), sau vài bước
đường đi tới một đỉnh \(B\) ở bên trái, rồi tới một đỉnh \(C\) ở bên phải.
Trong các sơ đồ dưới đây, cạnh liền biểu diễn một bước đi trực tiếp, còn cạnh
đứt biểu diễn một đoạn gồm nhiều bước đi.

\begin{center}
\begin{tikzpicture}[
  point/.style={circle,draw,fill=white,inner sep=0pt,minimum size=0.18cm},
  edge/.style={->,thick},
  multi/.style={->,thick,dashed},
  every node/.style={font=\small}
]
  \node[point,label=right:\(1\)] (p1) at (1.6,2.6) {};
  \node[point,label=below:\(A\)] (pa) at (1.15,0.1) {};
  \node[point,label=left:\(B\)] (pb) at (-0.35,1.25) {};
  \node[point,label=right:\(C\)] (pc) at (3.0,1.35) {};
  \draw[edge] (p1) -- (pa);
  \draw[multi] (pa) -- (pb);
  \draw[edge] (pb) -- (pc);
  \node[align=center] at (1.35,-0.65) {cấu hình có hai cạnh cắt nhau};

  \node[point,label=right:\(1\)] (q1) at (6.2,2.6) {};
  \node[point,label=below:\(A\)] (qa) at (5.75,0.1) {};
  \node[point,label=left:\(B\)] (qb) at (4.25,1.25) {};
  \node[point,label=right:\(C\)] (qc) at (7.6,1.35) {};
  \draw[edge] (q1) -- (qb);
  \draw[multi] (qb) -- (qa);
  \draw[edge] (qa) -- (qc);
  \node[align=center] at (5.95,-0.65) {đổi thứ tự để tránh cắt nhau};
\end{tikzpicture}
\end{center}

Ta có thể biến đổi thứ tự các đoạn đường để thay đoạn vượt qua hai phía bằng
một cách nối không cắt. Sau biến đổi, độ dài từ 1 tới đỉnh tiếp theo rõ ràng
ngắn hơn trước; do đó phương án trước biến đổi không thể là tối ưu. Vì thế,
từ 1, bước đầu tiên chỉ có thể tới 2 hoặc \(n\).

Tương tự, nếu bước đầu tới 2 thì bước thứ hai chỉ có thể tới \(n\) hoặc 3; nếu
bước đầu tới \(n\) thì bước thứ hai chỉ có thể tới \(n-1\) hoặc 2. Suy tiếp
như vậy, trong phương án tối ưu cuối cùng chắc chắn không có hai cạnh cắt
nhau. Phỏng đoán 2 được chứng minh.

Phần còn lại trở nên dễ xử lý. Theo chứng minh, xuất phát từ 1 thì bước đầu
chỉ có thể tới 2 hoặc \(n\). Nếu bước đầu tới 2, bài toán chuyển thành
``lấy 2 làm điểm xuất phát và đi qua mỗi đỉnh trong \(\{2,\ldots,n\}\) đúng
một lần''. Nếu bước đầu tới \(n\), bài toán chuyển thành ``lấy \(n\) làm điểm
xuất phát và đi qua mỗi đỉnh trong \(\{2,\ldots,n\}\) đúng một lần''.

Bài toán con sau chuyển đổi có cùng cấu trúc với bài toán ban đầu, chỉ nhỏ
hơn về kích thước, nên có thể dùng quy hoạch động. Gọi \(f[s,L,0]\) là độ dài
ngắn nhất khi xuất phát từ \(s\) và đi qua mỗi đỉnh trong
\(\{s,\ldots,s+L-1\}\) đúng một lần; gọi \(f[s,L,1]\) là độ dài ngắn nhất khi
xuất phát từ \(s+L-1\) và đi qua mỗi đỉnh trong \(\{s,\ldots,s+L-1\}\) đúng
một lần. Khi đó:
\[
\begin{aligned}
f[s,L,0]=\min\{&
  \operatorname{dis}(s,s+1)+f[s+1,L-1,0],\\
& \operatorname{dis}(s,s+L-1)+f[s+1,L-1,1]\},\\
f[s,L,1]=\min\{&
  \operatorname{dis}(s+L-1,s+L-2)+f[s,L-1,1],\\
& \operatorname{dis}(s+L-1,s)+f[s,L-1,0]\},\\
f[s,1,0]&=0,\qquad f[s,1,1]=0.
\end{aligned}
\]
Ở đây \(\operatorname{dis}(a,b)\) là khoảng cách Euclid giữa đỉnh thứ \(a\)
và đỉnh thứ \(b\). Độ phức tạp thời gian của thuật toán là \(O(n^2)\).

\paragraph{Nhận xét về bài toán.}
Ví dụ 1 là một bài tương đối khó. Điểm đột phá nằm ở việc tương tự với ví dụ
dẫn nhập: sự giống nhau về trực quan hình học kích hoạt một phỏng đoán mạnh,
giúp phân tích tính chất của đường tối ưu và tạo điều kiện cho quy hoạch
động. Ngay cả cách chứng minh phỏng đoán cũng liên quan chặt chẽ với một số
tư tưởng trong phần phân tích ví dụ dẫn nhập. Điều này không tách rời sự
giống nhau về bản chất và đặc điểm của hai bài toán.

\paragraph{Nhận xét của mục.}
Điều kiện tiên quyết để tương tự với bài toán quen thuộc là phân tích sự vật
một cách sâu sắc và thấu đáo. Chỉ khi phân tích được tính chất bản chất của
sự vật, ta mới có thể liên hệ các mô hình tương tự bị che khuất bởi vô số yếu
tố cụ thể phức tạp, và mới có thể dựa trên tình hình cụ thể của mô hình để
đưa ra phỏng đoán hợp lý.

Đồng thời, quan sát và liên tưởng là nền tảng để thực hiện tương tự. Nhờ quan
sát, ta thu được mô hình bản chất của bài toán; nhờ liên tưởng, ta liên hệ mô
hình ấy với bài toán quen thuộc và tương tự, rồi nêu phỏng đoán.

\subsection{Tương tự với bài toán đặc biệt}

Tương tự bằng đặc biệt hóa cũng là một loại tương tự thường dùng. Nó đặc biệt
hóa, cực đoan hóa hoặc biên giới hóa một số điều kiện trong bài toán; thông
qua phân tích bài toán ở tình huống cực đoan, ta cố gắng tìm quy luật giải
bài, rồi so sánh với bài toán ban đầu để nêu phỏng đoán. Có thể hình dung:
\[
\begin{array}{ccc}
\text{bài toán }A & \longrightarrow & \text{thuật toán }A\\
\big\uparrow\scriptstyle{\text{đặc biệt hóa}} && \big\downarrow\scriptstyle{\text{phỏng đoán}}\\
\text{bài toán }B & \longrightarrow & \text{thuật toán }B
\end{array}
\]

Các bước thông thường của tương tự đặc biệt là:
\begin{enumerate}
  \item Quan sát đặc trưng của điều kiện đề bài.
  \item Đặc biệt hóa vị trí, số lượng, trạng thái hoặc những điều kiện khác.
  \item Phân tích bài toán đặc biệt.
  \item So sánh bài toán đặc biệt với bài toán ban đầu.
  \item Dựa trên liên hệ và khác biệt giữa bài toán đặc biệt và bài toán ban
  đầu, lấy lời giải của bài toán đặc biệt làm khuôn mẫu để phỏng đoán cách
  giải bài toán ban đầu.
\end{enumerate}

\subsubsection*{Ví dụ 2: bài toán điểm và đường tròn}

\paragraph{Mô tả bài toán.}
Trên mặt phẳng có \(n\) điểm, trong đó \(3<n\le 10000\) và \(n\) là số lẻ.
Không có ba điểm thẳng hàng và không có bốn điểm đồng viên. Hãy lập trình tìm
ba điểm xác định một đường tròn sao cho \((n-3)/2\) điểm nằm ngoài đường
tròn và \((n-3)/2\) điểm nằm trong đường tròn.

\paragraph{Dữ liệu vào.}
Dòng đầu là một số nguyên \(n\). Tiếp theo là \(n\) dòng, mỗi dòng có hai số
nguyên \(x,y\), biểu diễn tọa độ đỉnh thứ \(i\).

\paragraph{Dữ liệu ra.}
Tệp ra chỉ có một dòng gồm ba số nguyên, lần lượt là chỉ số ba điểm được
chọn.

\paragraph{Phân tích.}
Sau khi đọc đề, phản xạ tự nhiên là liệt kê: chọn ba điểm xác định một đường
tròn, rồi lần lượt xét vị trí của các điểm còn lại so với đường tròn cho đến
khi tìm được nghiệm. Nhưng vết thương chí mạng của thuật toán này là độ phức
tạp thời gian quá cao, trong khi \(n\le 10000\). Ta phải tìm cách khác.

Vì quan hệ giữa \(n\) điểm rời rạc và phân tán gây khó cho suy nghĩ, ta xét
việc đặc biệt hóa vị trí của chúng: giả sử \(n\) điểm tạo thành một bao lồi.
Đồng thời đặt \(n=5\), là giá trị nhỏ nhất.

Trong tình huống đặc biệt này, giả sử năm điểm \(A,B,C,D,E\) tạo thành một
bao lồi. Chọn ba điểm \(A,C,D\) để xác định đường tròn thì \(B\) nằm ngoài
đường tròn, còn \(E\) nằm trong đường tròn.

Vì sao chọn \(A,C,D\) có thể khiến \(B\) và \(E\) nằm hai phía trong, ngoài
đường tròn? Vì năm điểm tạo thành một bao lồi, nên \(A,B,E\) đều ở cùng phía
của \(CD\). Do đó \(\angle CAD,\angle CBD,\angle CED\) đều là các góc chắn
cùng phía của \(CD\), và
\[
  \angle CBD<\angle CAD<\angle CED.
\]
Vì \(A\) nằm trên đường tròn, \(B\) nằm ngoài và \(E\) nằm trong.

Từ trường hợp đặc biệt này, ta khái quát được hai yếu tố:
\begin{enumerate}
  \item \(A,B,E\) nằm cùng phía của \(CD\).
  \item \(\angle CBD<\angle CAD<\angle CED\).
\end{enumerate}
Theo yếu tố 1, ba góc \(\angle CAD,\angle CBD,\angle CED\) đều là góc chắn
cùng phía của \(CD\), nên có thể dùng chúng để xác định quan hệ của \(A,B,E\)
với đường tròn. Yếu tố 2 bảo đảm \(B\) và \(E\) nằm lần lượt ngoài và trong
đường tròn. Hai yếu tố này là bảo đảm căn bản cho tính đúng đắn của thuật
toán, thiếu một trong hai đều không được.

Tiếp theo, ta so sánh với bài toán ban đầu và nêu phỏng đoán tương tự với hai
yếu tố trên:

\begin{quote}
Nếu trong bài toán ban đầu ba điểm \(A,B,C\) thỏa:
\begin{enumerate}
  \item đường thẳng \(AB\) làm cho mọi điểm khác ngoài \(A,B\) đều nằm cùng
  một phía của \(AB\);
  \item sau khi sắp xếp mọi góc \(\angle AXB\), với \(X\ne A,B\), góc
  \(\angle ACB\) đứng ở vị trí thứ \((n-3)/2+1\);
\end{enumerate}
thì \(A,B,C\) chắc chắn là ba điểm thỏa yêu cầu.
\end{quote}

\paragraph{Chứng minh.}
Do điều kiện 1 được thỏa, mọi \(\angle AXB\) đều là góc chắn cùng phía của
\(AB\). Vì vậy nếu \(\angle AXB>\angle ACB\) thì \(X\) nằm trong đường tròn
\(ABC\); nếu \(\angle AXB<\angle ACB\) thì \(X\) nằm ngoài đường tròn
\(ABC\).

Góc \(\angle ACB\) đứng ở vị trí thứ \((n-3)/2+1\), nên có \((n-3)/2\) góc
lớn hơn nó và \((n-3)/2\) góc nhỏ hơn nó. Vì vậy có \((n-3)/2\) điểm nằm
ngoài đường tròn và \((n-3)/2\) điểm nằm trong đường tròn. Do đó \(A,B,C\)
là ba điểm cần tìm.

Ta có thể trước hết tìm bao lồi của \(n\) điểm, chọn tùy ý hai điểm kề nhau
\(A,B\) trên bao lồi; các điểm còn lại chắc chắn ở cùng phía của \(AB\). Sau
đó sắp xếp các điểm còn lại theo góc chắn với \(AB\), chọn điểm \(C\) đứng ở
vị trí thứ \((n-3)/2+1\). Khi đó \(A,B,C\) chính là đáp án.

\paragraph{Nhận xét về bài toán.}
Đột phá thực chất của bài này nằm ở việc đặc biệt hóa bài toán. Sau khi đặc
biệt hóa, quan hệ giữa các điểm trở nên ngắn gọn và rõ ràng, nên chỉ cần phân
tích nhẹ là giải được. Sau đó ta phân tích một số yếu tố giải bài trong tình
huống đặc biệt, rút ra một kết luận tổng quát hơn, rồi dựa vào đó phỏng đoán
thuật toán cho bài toán ban đầu.

\paragraph{Nhận xét của mục.}
Tương tự với bài toán đặc biệt là một trong những phương pháp quan trọng để
có được phỏng đoán. Vì bài toán sau đặc biệt hóa vẫn giữ lại phần lớn đặc
trưng và tính chất của bài toán ban đầu, nên dùng kết quả nghiên cứu của nó
để phỏng đoán thường dễ hơn các loại tương tự khác.

Khi dùng tương tự bằng đặc biệt hóa, cần đặc biệt chú ý so sánh điểm khác
giữa bài toán đặc biệt và bài toán ban đầu, tuyệt đối không mở rộng thuật
toán đặc biệt một cách máy móc.

Điều kiện tiên quyết của tương tự đặc biệt là tích lũy. Nên đặc biệt hóa điều
kiện nào? Đặc biệt hóa thành dạng nào? Những điều này đều phụ thuộc vào việc
thường ngày rèn luyện, để ý và tích lũy kinh nghiệm giải bài. Không có tích
lũy vững chắc thì ``tương tự đặc biệt'' chỉ là một lâu đài trên không.

\subsection{Nhận xét về tương tự}

Tương tự là một phương pháp quan trọng để đi tới phỏng đoán. Tương tự với bài
toán đặc biệt và tương tự với bài toán quen thuộc là hai hình thức cơ bản của
tương tự.

Hình thức thứ nhất so sánh bài toán xa lạ với một bài toán quen thuộc tương
tự, rồi phỏng đoán cách giải bài toán xa lạ. Đó là đi từ xa lạ tới quen
thuộc.

Hình thức thứ hai làm bài toán phức tạp trở nên đơn giản và đặc biệt, thông
qua phân tích bài toán đơn giản để rút ra quy luật giải bài nói chung, rồi so
sánh với bài toán ban đầu để giải quyết vấn đề. Đó là đi từ tổng quát tới đặc
biệt, từ phức tạp tới đơn giản.

Bản chất của hai phương pháp đều là ``quy về'' những điều đã biết, tức chuyển
từ thế giới đã biết sang thế giới chưa biết. Việc chuyển đổi ấy phải xây dựng
trên tích lũy, quan sát, phân tích và liên tưởng, chứ tuyệt đối không chỉ là
may mắn.

\section{Phỏng đoán bằng quy nạp}

\subsection{Cơ sở lý thuyết của phỏng đoán quy nạp}

Quy nạp là một hình thức suy luận rút ra kết luận phổ quát thông qua phân
tích các trường hợp riêng. Nó gồm hai phần: tiền đề và kết luận. Tiền đề là
một số sự thật riêng lẻ đã biết, tức các phán đoán hoặc phát biểu cá biệt hay
đặc biệt; kết luận là phỏng đoán thu được bằng suy luận từ tiền đề, tức một
phát biểu hoặc phán đoán phổ quát. Mô thức tư duy là: giả sử \(M_i\)
\((i=1,2,\ldots,n)\) là các trường hợp riêng hoặc tập con của đối tượng cần
nghiên cứu \(M\). Nếu các \(M_i\) đều có tính chất \(P\), thì từ đó phỏng
đoán rằng \(M\) cũng có thể có tính chất \(P\).

Nếu
\[
  \bigcup_{i=1}^{n} M_i=M,
\]
thì phương pháp quy nạp này gọi là \term{quy nạp đầy đủ}. Vì nó xét hết mọi
trường hợp riêng của đối tượng nghiên cứu, kết luận là đúng và đáng tin cậy.
Quy nạp đầy đủ có thể dùng như một phương pháp chứng minh, còn gọi là quy nạp
bằng liệt kê.

Nếu
\[
  \bigcup_{i=1}^{n} M_i
\]
là một tập con thật sự của \(M\), thì phương pháp này gọi là \term{quy nạp
không đầy đủ}. Vì quy nạp không đầy đủ không xét hết mọi đối tượng nghiên
cứu, kết luận thu được chỉ có thể xem là phỏng đoán; nó đúng hay sai còn phải
được chứng minh thêm hoặc bác bỏ bằng phản ví dụ.

Trong tin học, ta thường trước hết dùng ``quy nạp không đầy đủ'' để thu được
kết luận phổ quát, rồi dùng ``quy nạp đầy đủ'' để chứng minh nó, thường là
bằng quy nạp toán học.

Các bước thông thường của phỏng đoán quy nạp là:
\begin{enumerate}
  \item \term{Liệt kê.} Liệt kê một số dữ liệu đặc biệt, đơn giản, quy mô
  nhỏ. Bước này có thể làm tay hoặc viết một chương trình tìm kiếm nhỏ.
  \item \term{Quan sát.} Quan sát quy luật của dữ liệu đã liệt kê.
  \item \term{Phỏng đoán.} Dựa vào một phần dữ liệu để phỏng đoán kết luận
  tổng quát.
  \item \term{Chứng minh.} Chứng minh tính đúng đắn của phỏng đoán, thường
  bằng quy nạp toán học.
\end{enumerate}

\subsection{Ứng dụng của phỏng đoán quy nạp trong tin học}

\subsubsection*{Ví dụ 3: ếch qua sông}

\paragraph{Mô tả bài toán.}
Một đội ếch có kích thước đôi một khác nhau đang đứng trên trụ đá ở bờ trái,
ký hiệu là \(A\), và muốn sang trụ đá ở bờ phải, ký hiệu là \(D\). Giữa sông
có vài lá sen \(Y_1,\ldots,Y_m\) và vài trụ đá \(S_1,\ldots,S_n\).

Quy tắc xếp hàng và di chuyển của ếch như sau:
\begin{enumerate}
  \item Mỗi con ếch chỉ được đứng trên lá sen, trụ đá, hoặc trên lưng con ếch
  lớn hơn nó đúng một cỡ. Các vị trí này gọi chung là điểm đặt chân hợp lệ.
  \item Một con ếch chỉ có thể rời khỏi một điểm đặt chân để tới một điểm đặt
  chân khác khi không có con ếch nào đứng trên lưng nó.
  \item Ếch được phép nhảy trực tiếp từ bờ trái \(A\) tới trụ đá hoặc lá sen
  giữa sông, cũng như tới trụ đá bờ phải \(D\). Ếch cũng được phép nhảy từ
  trụ đá hoặc lá sen giữa sông tới \(D\).
  \item Ếch có thể nhảy qua lại giữa các trụ đá, giữa các lá sen, và giữa trụ
  đá với lá sen ở giữa sông.
  \item Sau khi rời trụ đá bờ trái, ếch không được quay lại bờ trái; sau khi
  tới bờ phải, nó không được nhảy ngược trở lại.
  \item Giả sử sức chịu tải của trụ đá rất lớn, cho phép bao nhiêu con ếch
  đứng trực tiếp lên cũng được. Các con ếch khác chỉ được đứng theo quy tắc
  1, tức trên lưng con lớn hơn nó đúng một cỡ.
  \item Lá sen không chỉ nhỏ về diện tích mà sức chịu tải cũng hạn chế, nhiều
  nhất chỉ có một con ếch đứng trên đó.
  \item Mỗi bước chỉ được di chuyển một con ếch, và sau khi di chuyển phải
  thỏa quy tắc xếp hàng.
  \item Ban đầu tất cả ếch đều đứng trên \(A\). Con lớn nhất đứng trực tiếp
  trên trụ đá, còn các con khác theo quy tắc 6 đứng trên lưng con lớn hơn nó
  đúng một cỡ.
\end{enumerate}
Ếch muốn cuối cùng toàn bộ đội chuyển tới \(D\) và hoàn thành việc xếp hàng.
Giả sử giữa sông có \(m\) lá sen và \(n\) trụ đá. Hãy tìm số ếch nhiều nhất
có thể đi từ \(A\) sang \(D\) mà vẫn thỏa mọi quy tắc.

\paragraph{Dữ liệu vào.}
Tệp vào chỉ có hai dòng, mỗi dòng chứa một số nguyên. Dòng đầu là số trụ đá
giữa sông \(n\), với \(0\le n\le 25\). Dòng thứ hai là số lá sen \(m\), với
\(0\le m\le 25\).

\paragraph{Dữ liệu ra.}
Tệp ra chỉ chứa một số nguyên: số ếch nhiều nhất có thể qua sông.

\paragraph{Phân tích.}
Thoạt nhìn, bài này không có manh mối rõ ràng. Chú ý rằng kết quả cuối cùng
chỉ phụ thuộc vào hai tham số \(n,m\), nên ta bước đầu phỏng đoán rằng kết
quả là một biểu thức toán học theo \(n,m\). Bắt đầu từ các trường hợp đơn
giản để xem quy luật.

Gọi \(f[n,m]\) là số ếch có thể từ bờ trái nhảy sang bờ phải nhờ \(n\) trụ đá
và \(m\) lá sen. Khi \(n=0\), rõ ràng \(f[0,m]=m+1\). Khi \(m=0\), viết một
chương trình tìm kiếm đơn giản cho một số kết quả:
\[
  f[1,0]=2,\qquad f[2,0]=4,\qquad f[3,0]=8.
\]
Từ đó phỏng đoán \(f[n,0]=2^n\). Khi \(m=1\), theo kết quả tìm kiếm đơn giản:
\[
  f[1,1]=4,\qquad f[2,1]=8,\qquad f[3,1]=16,
\]
nên phỏng đoán \(f[n,1]=2^{n+1}\).

Tổng hợp lại, đều dưới dạng phỏng đoán:
\[
\begin{aligned}
m=0&:\quad f[n,0]=2^n,\\
m=1&:\quad f[n,1]=2^{n+1},\\
n=0&:\quad f[0,m]=m+1.
\end{aligned}
\]
Quan sát các công thức ấy, ta liên tưởng:
\[
\begin{aligned}
m=0&:\quad f[n,0]=2^n=(0+1)2^n=(m+1)2^n,\\
m=1&:\quad f[n,1]=2^{n+1}=2\cdot 2^n=(m+1)2^n,\\
n=0&:\quad f[0,m]=m+1=(m+1)\cdot 1=(m+1)2^n.
\end{aligned}
\]
Vì thế ta nêu phỏng đoán:
\[
  f[n,m]=(m+1)2^n.
\]

Tiếp theo chứng minh phỏng đoán này.

Khi \(n=0\), trước hết cho \(m\) con ếch nhảy lên các lá sen; sau đó cho con
ếch số \(m+1\) nhảy tới \(D\); cuối cùng cho \(m\) con trên lá sen lần lượt
nhảy tới \(D\) theo thứ tự số hiệu từ lớn đến nhỏ. Như vậy có thể cho tổng
cộng \(m+1=(m+1)2^0\) con ếch qua sông, tức \(f[0,m]=m+1\).

Khi \(n\ne 0\), hiển nhiên trước hết phải cho các con ếch có số hiệu nhỏ hơn
\(f[n,m]\) nhảy tới lá sen hoặc trụ đá để tạm đứng; sau đó cho con ếch số
\(f[n,m]\) nhảy tới \(D\); cuối cùng cho \(f[n,m]-1\) con ếch đang ở lá sen
và trụ đá nhảy tới \(D\). Vì vậy để số ếch qua sông là lớn nhất, ta phải làm
sao cho càng nhiều ếch càng tốt tạm đứng trên lá sen và trụ đá.

Trước hết xét việc làm cho trụ đá số \(n\) chứa càng nhiều ếch càng tốt. Có
thể sử dụng \(n-1\) trụ đá còn lại và \(m\) lá sen, nên trước hết có thể cho
\(f[n-1,m]\) con ếch nhảy lên trụ đá số \(n\). Tương tự, nhiều nhất có thể
cho \(f[n-2,m]\) con ếch nhảy lên trụ đá số \(n-1\), và cứ thế, nhiều nhất
có thể cho \(f[0,m]\) con ếch nhảy lên trụ đá số 1.

Do đó
\[
  f[n,m]=f[n-1,m]+f[n-2,m]+\cdots+f[0,m]+m+1.
\]
Có thể thấy tham số \(m\) ở đây đóng vai trò hằng. Đặt \(g[n]=f[n,m]\), ta
có:
\[
\begin{aligned}
g[i]&=g[0]+g[1]+\cdots+g[i-1]+m+1,\\
g[0]&=m+1.
\end{aligned}
\]
Cần tìm \(g[n]\). Phân tích đơn giản cho thấy
\[
  f[n,m]=g[n]=(m+1)2^n.
\]
Vì vậy phỏng đoán là đúng.

\paragraph{Nhận xét về bài toán.}
Đây là bài thứ hai của ngày thi thứ hai NOI 2000, có độ khó nhất định. Nếu
không dùng phỏng đoán quy nạp thì rất khó làm; ngay cả khi làm được cũng sẽ
tốn thời gian đáng kể.

Cần chú ý rằng phỏng đoán thu được bằng quy nạp chưa chắc đúng, mà phải được
chứng minh nghiêm ngặt. Phương pháp chứng minh phần lớn là quy nạp toán học;
điều này do chính đặc điểm của phỏng đoán quy nạp quyết định.

\subsubsection*{Ví dụ 4: bài toán hoán vị}

\paragraph{Mô tả bài toán.}
Trong các hoán vị của các số nguyên \(1,2,\ldots,N\), có một số hoán vị thỏa
tính chất \(A\): ngoài số nguyên cuối cùng, sau mỗi số nguyên trong hoán vị
đều có, không nhất thiết ngay sát, một số nguyên khác chênh lệch với nó đúng
1. Ví dụ khi \(N=4\), hoán vị \(1432\) có tính chất \(A\), còn \(2431\) thì
không.

Cho một hoán vị gồm \(N\) số, trong đó đã biết giá trị ở \(P\) vị trí
\((P\le N)\). Hỏi có bao nhiêu hoán vị vừa có các giá trị đã biết ở \(P\) vị
trí ấy, vừa có tính chất \(A\). Ví dụ khi \(N=4\), biết vị trí thứ nhất và
thứ hai lần lượt là 1 và 4, thì \(1432\) và \(1423\) là hai hoán vị như vậy.

\paragraph{Dữ liệu vào.}
Dòng đầu là \(N\), với \(N\le 50\). Dòng thứ hai có \(N\) số, trong đó số
chưa biết được biểu diễn bằng 0.

\paragraph{Dữ liệu ra.}
In ra tổng số hoán vị thỏa các điều kiện trên.

\paragraph{Ví dụ.}
\begin{verbatim}
Input
4
1400

Output
2
\end{verbatim}

\paragraph{Phân tích.}
Bài toán có tính che giấu khá mạnh. Điều kiện mà dãy cần thỏa tương đối phức
tạp, khiến ta khó bắt đầu. Ta thử viết ra một số dãy và quan sát quy luật.
Khi \(n=5\), có 16 dãy thỏa yêu cầu:
\[
\begin{array}{llll}
12345 & 12354 & 12534 & 12543\\
15234 & 15243 & 15423 & 15432\\
51234 & 51243 & 51423 & 51432\\
54123 & 54132 & 54312 & 54321
\end{array}
\]
Quan sát kỹ ta thấy một quy luật: với bất kỳ hoán vị nào, \(k\) vị trí cuối
\((1\le k\le n)\) đều tạo thành một tập gồm các số nguyên liên tiếp. Ví dụ
với \(51234\), một vị trí cuối là \(\{4\}\), hai vị trí cuối là \(\{3,4\}\),
ba vị trí cuối là \(\{2,3,4\}\), bốn vị trí cuối là \(\{1,2,3,4\}\), năm vị
trí cuối là \(\{1,2,3,4,5\}\). Các hoán vị còn lại cũng thỏa quy luật này.

Từ đó ta có thể mạnh dạn nêu:

\begin{quote}
\term{Phỏng đoán 1.} Với bất kỳ hoán vị thỏa yêu cầu nào, \(k\) vị trí cuối
\((1\le k\le n)\) đều tạo thành một tập gồm các số nguyên liên tiếp.
\end{quote}

\paragraph{Chứng minh.}
Ký hiệu phần tử ở vị trí thứ \(i\) của dãy là \(v[i]\). Giả sử \(x\) vị trí
cuối của dãy, với \(x\ge 1\), là một số các số nguyên liên tiếp tạo thành tập
\(\{s,\ldots,t\}\). Khi đó số ở vị trí ngược thứ \(x+1\), tức
\(v[n-x+1]\), chắc chắn bằng \(p-1\) hoặc \(p+1\) với một
\(p\in\{s,\ldots,t\}\). Rõ ràng
\[
  \{v[n-x+1]\}\cup\{s,\ldots,t\}
\]
bằng \(\{s-1,\ldots,t\}\) hoặc \(\{s,\ldots,t+1\}\), vẫn là một tập gồm các
số nguyên liên tiếp.

Dựa trên chứng minh này, ta phỏng đoán tiếp:

\begin{quote}
\term{Phỏng đoán 2.} Nếu với một hoán vị, \(k\) vị trí cuối
\((1\le k\le n)\) đều tạo thành một tập gồm các số nguyên liên tiếp, thì hoán
vị ấy thỏa yêu cầu của đề.
\end{quote}

\paragraph{Chứng minh.}
Ký hiệu phần tử ở vị trí thứ \(i\) là \(v[i]\). Giả sử
\[
  \{v[n],v[n-1],\ldots,v[n-k+1]\}=\{s,\ldots,t\}.
\]
Vì \(\{v[n-1],\ldots,v[n-k+1]\}\) cũng là một tập gồm các số nguyên liên
tiếp, nên \(v[n]\) bằng \(s\) hoặc \(t\). Nếu \(v[n]=s\), thì chắc chắn
\[
  s+1\in\{v[n-1],\ldots,v[n-k+1]\}.
\]
Nếu \(v[n]=t\), thì chắc chắn
\[
  t-1\in\{v[n-1],\ldots,v[n-k+1]\}.
\]
Nói cách khác, đây là một dãy thỏa yêu cầu.

Vì vậy bài toán chuyển thành: đếm số hoán vị của \(1,2,\ldots,N\) sao cho
\(k\) vị trí cuối \((1\le k\le n)\) đều tạo thành một tập gồm các số nguyên
liên tiếp.

Phần còn lại thuận lợi hơn. Theo hai phỏng đoán, ta dễ nghĩ tới quy hoạch
động. Gọi \(g[s,r]\) là số dãy \(C\) thỏa các điều kiện sau:
\begin{itemize}
  \item \(C\) được tạo bởi các số trong tập \(\{s,\ldots,s+r-1\}\);
  \item \(k\) vị trí cuối \((1\le k\le r)\) của \(C\) đều tạo thành một tập
  gồm các số nguyên liên tiếp;
  \item nếu trong bài toán ban đầu số ở vị trí ngược thứ \(i\) đã được xác
  định là \(x\), với \(1\le i\le r\), thì số ở vị trí ngược thứ \(i\) của
  \(C\) cũng phải là \(x\).
\end{itemize}
Khi đó:
\[
g[s,r]=
\begin{cases}
g[s+1,r-1], & \text{nếu vị trí ngược thứ } r \text{ được xác định là } s,\\
g[s,r-1], & \text{nếu vị trí ngược thứ } r \text{ được xác định là } s+r-1,\\
g[s,r-1]+g[s+1,r-1], & \text{nếu vị trí ngược thứ } r \text{ chưa xác định},\\
0, & \text{trong các trường hợp khác}.
\end{cases}
\]
Điều kiện đầu là
\[
  g[s,1]=1.
\]
Cần tìm \(g[1,n]\). Độ phức tạp thời gian của thuật toán là \(O(n^2)\).

\paragraph{Nhận xét về bài toán.}
Bài này là đề tập huấn đội tuyển tỉnh Hồ Nam, có độ khó khá cao. Nguyên nhân
chủ yếu là các điều kiện cho trước khá phức tạp và khó chuyển hóa để sử dụng.
Vì thế trong phòng thi, nó dễ khiến đa số thí sinh dùng tìm kiếm; khi đó chỉ
có hai thí sinh dùng quy hoạch động nhanh hơn.

Điểm đột phá của bài nằm ở hai phỏng đoán rất táo bạo. Nhờ hai phỏng đoán
này, điều kiện khó dùng trực tiếp được chuyển thành một hình thức quen thuộc
và dễ sử dụng hơn, mở đường cho quy hoạch động.

Tuy nhiên từ ``liệt kê dữ liệu đơn giản'' tới ``đề xuất phỏng đoán'' vẫn là
một quãng đường đáng kể. Nếu không có tích lũy kinh nghiệm phong phú, linh
cảm giải bài, năng lực quan sát, khái quát và liên tưởng, rất khó phát hiện
được kết luận có giá trị.

\subsection{Nhận xét về quy nạp}

Mục này đã giới thiệu phân loại cơ bản, phương pháp cơ bản của quy nạp, cũng
như các ví dụ ứng dụng trong tin học.

Ngoài tương tự, quy nạp là một phương pháp quan trọng khác để đi tới phỏng
đoán. Tư tưởng cơ bản của nó là đi từ đặc biệt tới tổng quát, từ đơn giản tới
phức tạp, từ quy mô nhỏ tới quy mô lớn. Đây cũng là một tư tưởng ``quy về''
những điều đã biết.

Một điểm đáng chú ý nữa là quy nạp không nhất thiết trực tiếp phỏng đoán ra
kết quả của bài toán. Có thể như ví dụ 4: phỏng đoán nhận được bằng quy nạp
không gắn trực tiếp với kết quả cuối cùng, nhưng kết luận của phỏng đoán lại
giúp ích lớn cho hướng tư duy và hướng giải bài.

Khi tiến hành phỏng đoán quy nạp, cần chú ý tư duy phân kỳ và tưởng tượng đầy
đủ: nghĩ và liên hệ theo mọi hướng có thể. Đừng để cảm hứng lóe lên rồi trôi
mất.

\section{Kết luận}

Thật ra khi giải bài, ta đều vô thức phỏng đoán: phỏng đoán tính chất của bài
toán, phỏng đoán hình thức của mô hình, phỏng đoán nội dung của thuật toán.
Phỏng đoán là con đường quan trọng để thăm dò từ điều đã biết tới điều chưa
biết.

Phỏng đoán không đồng nghĩa với ngẫu nhiên cộng may mắn. Phỏng đoán là một
suy luận hợp lý nhưng chủ quan và chưa đầy đủ, xây dựng trên cơ sở quan sát,
phân tích, liên tưởng và quy nạp. Nó là biểu hiện tổng hợp của năng lực quan
sát, khái quát, liên tưởng và sáng tạo. Trình độ phỏng đoán phản ánh rõ nhất
chất lượng tư duy của một người.

Theo cách thu được phỏng đoán, có thể chia thành phỏng đoán bằng tương tự và
phỏng đoán bằng quy nạp. Các phần trên đã trình bày chi tiết hai loại này.

Nhưng dù dùng phương pháp nào, phỏng đoán cũng không thể tách khỏi quan sát,
phân tích và thực hành. Có thể nói: quan sát là dòng máu của phỏng đoán,
xuyên suốt toàn bộ quá trình; phân tích là linh hồn của phỏng đoán, là tiền
đề và nền tảng để nêu phỏng đoán; thực hành là bộ khung của phỏng đoán, là
nền móng nâng đỡ phỏng đoán. Rời khỏi thực tế thì phỏng đoán không có ý
nghĩa.

Khi đề thi tin học phát triển theo hướng che giấu hơn, đa chiều hơn và trừu
tượng hơn, độ khó của đề ngày càng tăng, yêu cầu đối với năng lực sáng tạo
của thí sinh cũng ngày càng cao. Chiến lược giải bài quan trọng là phỏng
đoán chắc chắn sẽ giữ vai trò ngày càng lớn.

\section*{Tài liệu tham khảo}

\begin{enumerate}
  \item Ouyang Weicheng, \emph{Nghiên cứu phương pháp giải toán sơ cấp}, Nhà
  xuất bản Giáo dục Hồ Nam.
  \item Đề tập huấn Hồ Nam năm 2000.
  \item Đề thi NOI 2000.
  \item Tài liệu liên quan trên Mạng giáo dục cơ sở Trung Quốc,
  \texttt{www.cbe21.com}.
\end{enumerate}

\appendix
\section{Bản trình chiếu kèm theo trong PDF}

Phần sau dịch nội dung bản trình chiếu xuất hiện sau bài viết trong cùng tệp
PDF. Nội dung phần lớn trùng với bài viết chính, nhưng được giữ lại vì văn
bản có thể trích xuất được.

\subsection{Trang tiêu đề}

\begin{center}
\Large Phỏng đoán và ứng dụng của nó
\end{center}

\subsection{Nội dung}

\begin{itemize}
  \item Lời mở đầu.
  \item Phỏng đoán bằng tương tự.
  \item Tương tự với bài toán quen thuộc: người khổng lồ và ma
  \(\leftrightarrow\) nỗi phiền của ếch.
  \item Tương tự với bài toán đặc biệt: bài toán điểm và đường tròn.
  \item Phỏng đoán bằng quy nạp: ếch qua sông, bài toán hoán vị.
  \item Lời kết.
\end{itemize}

\subsection{Lời mở đầu}

Phỏng đoán là một suy luận hợp lý nhưng chủ quan và chưa đầy đủ. Nó được xây
dựng trên cơ sở quan sát, phân tích, liên tưởng và quy nạp.

Hai loại chính là phỏng đoán bằng tương tự và phỏng đoán bằng quy nạp. Quá
trình cơ bản là:
\[
\begin{array}{ccccc}
\text{đề xuất phỏng đoán}
&\longrightarrow&
\text{kiểm tra bằng dữ liệu nhỏ}
&\xrightarrow{\text{đúng}}&
\text{chứng minh kết luận}\\
&&\bigg\downarrow\scriptstyle{\text{sai}}&&\bigg\downarrow\\
&&\text{sửa phỏng đoán}
&\longleftarrow&
\text{giải quyết vấn đề}
\end{array}
\]
Do đó, phỏng đoán là quá trình dựa trên phân tích sâu, không ngừng thăm dò và
liên tục sửa chữa; nó không thể hoàn thành tức thời. Đề xuất phỏng đoán là
phần lõi và trọng điểm của toàn bộ quá trình.

\subsection{Phỏng đoán bằng tương tự}

\[
\begin{array}{ccc}
\text{sự vật }A && \text{sự vật }B\\
\text{tính chất }a && \text{tính chất }a'\\
\text{tính chất }b && \text{tính chất }b'\\
\text{tính chất }c && \text{tính chất }c'\\
\text{tính chất }d & \xrightarrow{\text{tương tự, phỏng đoán}} &
\text{tính chất }d'
\end{array}
\]

\subsection{Ví dụ dẫn nhập: người khổng lồ và ma}

Trên mặt phẳng có \(n\) người khổng lồ và \(n\) con ma
\((1\le n\le 50)\). Mỗi người khổng lồ có vũ khí laser kiểu mới để tấn công
ma theo đường thẳng, nhưng hai tia laser không được cắt nhau. Không có ba cá
thể bất kỳ cùng nằm trên một đường thẳng. Mỗi người khổng lồ tấn công một con
ma. Hãy tìm một phương án tấn công bất kỳ.

Bài toán thực chất là xác định song ánh giữa người khổng lồ và ma, nên gợi
tới ghép cặp. Xây dựng đồ thị hai phía với tập đỉnh \(X\) là người khổng lồ,
tập đỉnh \(Y\) là ma, và nối mọi cạnh \(x\in X,y\in Y\). Nhiệm vụ là tìm một
ghép hoàn hảo có các cạnh không cắt nhau.

Giả sử có hai cạnh cắt nhau. Khi đổi cặp nối, ta có
\[
  OA+OB>AB,\qquad OC+OD>CD,\qquad AD+CB>AB+CD.
\]
Vì vậy phương án mới có tổng độ dài cạnh ghép nhỏ hơn. Do đó mọi ghép hoàn
hảo chứa cạnh cắt nhau đều có thể chuyển thành ghép hoàn hảo có tổng độ dài
nhỏ hơn; ghép hoàn hảo có tổng độ dài nhỏ nhất chắc chắn không chứa cạnh cắt
nhau. Ta lập mô hình ghép tối ưu trên đồ thị hai phía, đặt trọng số cạnh là
khoảng cách Euclid, rồi dùng thuật toán Hungary để tìm ghép tối ưu.

Nhận xét: mỗi đỉnh trong đồ thị tương ứng với một điểm trên mặt phẳng, nên
giữa các đỉnh có cả quan hệ logic lẫn quan hệ trực quan hình học. Tận dụng
quan hệ ẩn này giúp mô hình thuật toán sáng rõ.

\subsection{Ví dụ 1: nỗi phiền của ếch}

Trong ao có \(n\) lá sen \((1\le n\le 1000)\), đúng bằng các đỉnh của một đa
giác lồi và được đánh số theo chiều kim đồng hồ \(1,2,\ldots,n\). Một chú ếch
đứng ở lá số 1, muốn đi qua mỗi lá đúng một lần và muốn tổng quãng đường
ngắn nhất. Mô hình là đường Hamilton ngắn nhất xuất phát từ 1 trong đồ thị
đầy đủ có trọng số Euclid.

Quan sát đồ thị:
\begin{itemize}
  \item đó là đồ thị đầy đủ;
  \item trọng số cạnh là khoảng cách Euclid giữa hai đầu mút;
  \item mỗi đỉnh tương ứng với một điểm trên mặt phẳng.
\end{itemize}
Các tính chất này rất giống ví dụ dẫn nhập. Nhờ đó ta phỏng đoán:

\begin{quote}
Đường Hamilton ngắn nhất không có hai cạnh cắt nhau.
\end{quote}

Nếu xuất phát từ 1 và đi đầu tiên tới \(A\), với \(2<A<n\), thì các đỉnh còn
lại nằm hai phía của đoạn \(\langle 1,A\rangle\). Nếu đường đi từ \(A\) sau
vài bước tới một đỉnh \(B\) ở một phía rồi tới \(C\) ở phía kia, ta có thể đổi
thứ tự các đoạn để được đường ngắn hơn. Vì thế phương án trước không tối ưu,
và đường Hamilton tối ưu không chứa cạnh cắt nhau.

Kết quả là từ 1, bước đầu chỉ tới 2 hoặc \(n\). Nếu tới 2, bài toán con là
xuất phát từ 2 và đi qua \(\{2,\ldots,n\}\); nếu tới \(n\), bài toán con là
xuất phát từ \(n\) và đi qua \(\{2,\ldots,n\}\). Đặt:
\[
\begin{aligned}
f[s,L,0]&=\text{độ dài ngắn nhất khi xuất phát từ }s
\text{ và đi qua }\{s,\ldots,s+L-1\},\\
f[s,L,1]&=\text{độ dài ngắn nhất khi xuất phát từ }s+L-1
\text{ và đi qua }\{s,\ldots,s+L-1\}.
\end{aligned}
\]
Khi đó:
\[
\begin{aligned}
f[s,L,0]=\min\{&
\operatorname{dis}(s,s+1)+f[s+1,L-1,0],\\
&\operatorname{dis}(s,s+L-1)+f[s+1,L-1,1]\},\\
f[s,L,1]=\min\{&
\operatorname{dis}(s+L-1,s+L-2)+f[s,L-1,1],\\
&\operatorname{dis}(s+L-1,s)+f[s,L-1,0]\},\\
f[s,1,0]&=f[s,1,1]=0.
\end{aligned}
\]
Cần tìm \(f[1,n,0]\).

\subsection{Nhận xét về tương tự với bài quen thuộc}

\[
\begin{array}{ccc}
\text{bài toán }A & \xrightarrow{\text{tương tự}} & \text{bài toán }B\\
\downarrow\scriptstyle{\text{giải}} && \downarrow\scriptstyle{\text{giải}}\\
\text{thuật toán }A & \xrightarrow{\text{phỏng đoán}} & \text{thuật toán }B
\end{array}
\]

Cần phân tích mô hình sâu sắc và thấu đáo; cần quan sát và liên tưởng. Các
bước là phân tích đặc trưng để trừu tượng hóa mô hình, liên tưởng tới bài
quen thuộc có mô hình tương tự, so sánh điểm giống và khác, rồi dựa vào bài
đã biết để phỏng đoán lời giải bài mới.

\subsection{Phỏng đoán bằng quy nạp}

Quy nạp là hình thức suy luận rút ra kết luận phổ quát thông qua phân tích
các trường hợp riêng.

\[
\begin{array}{ccc}
\text{một số sự thật riêng lẻ, đặc biệt}
&\xrightarrow{\text{quy nạp}}&
\text{phát biểu phổ quát, phỏng đoán}
\end{array}
\]

Đó là quá trình quan sát, phân tích, liên tưởng và khái quát; là suy luận hợp
lý nhưng chủ quan và chưa đầy đủ, nên phải qua chứng minh logic nghiêm ngặt.

Các bước của phỏng đoán quy nạp:
\begin{itemize}
  \item liệt kê dữ liệu đặc biệt, đơn giản, quy mô nhỏ, bằng tay hoặc bằng
  chương trình tìm kiếm nhỏ;
  \item quan sát quy luật của dữ liệu đã liệt kê;
  \item phỏng đoán kết luận tổng quát từ dữ liệu bộ phận;
  \item chứng minh tính đúng đắn của phỏng đoán, thường bằng quy nạp toán học.
\end{itemize}

\subsection{Ví dụ 4: bài toán hoán vị}

Trong các hoán vị của \(1,2,\ldots,N\), một số hoán vị có tính chất \(A\):
ngoài số cuối cùng, sau mỗi số đều có một số chênh lệch với nó đúng 1, không
nhất thiết đứng ngay sau. Ví dụ khi \(N=4\), \(1432\) thỏa tính chất \(A\),
còn \(2431\) không thỏa. Biết một số vị trí trong hoán vị, cần đếm số hoán vị
thỏa các giá trị đã biết và tính chất \(A\).

Khi \(n=5\), các dãy tìm được bằng tìm kiếm là:
\[
\begin{array}{llll}
12345 & 12354 & 12534 & 12543\\
15234 & 15243 & 15423 & 15432\\
51234 & 51243 & 51423 & 51432\\
54123 & 54132 & 54312 & 54321
\end{array}
\]
Quan sát thấy: với bất kỳ hoán vị nào, \(k\) vị trí cuối
\((1\le k\le n)\) đều là một tập gồm các số nguyên liên tiếp. Ví dụ với
\(51243\):
\[
\begin{array}{c|c}
\text{số vị trí cuối} & \text{tập nhận được}\\ \hline
5 & \{1,\ldots,5\}\\
4 & \{1,\ldots,4\}\\
3 & \{2,\ldots,4\}\\
2 & \{3,\ldots,4\}\\
1 & \{3\}
\end{array}
\]

Phỏng đoán 1: với bất kỳ hoán vị thỏa yêu cầu nào, \(k\) vị trí cuối
\((1\le k\le n)\) đều tạo thành một tập gồm các số nguyên liên tiếp. Đặt phần
tử thứ \(i\) là \(v[i]\). Nếu \(x\) vị trí cuối là tập \(\{s,\ldots,t\}\),
thì phần tử ở vị trí ngược thứ \(x+1\) phải bằng \(p-1\) hoặc \(p+1\) với
\(p\in\{s,\ldots,t\}\); do đó khi thêm nó vào ta vẫn được một khoảng số
nguyên liên tiếp.

Phỏng đoán 2: nếu một hoán vị có \(k\) vị trí cuối
\((1\le k\le n)\) luôn tạo thành một tập gồm các số nguyên liên tiếp, thì
hoán vị ấy thỏa yêu cầu. Nếu
\[
\{v[n],v[n-1],\ldots,v[n-k+1]\}=\{s,\ldots,t\},
\]
thì vì tập bỏ \(v[n]\) cũng liên tiếp, \(v[n]=s\) hoặc \(v[n]=t\). Khi
\(v[n]=s\), trong phần sau có \(s+1\); khi \(v[n]=t\), trong phần sau có
\(t-1\). Vì vậy dãy thỏa tính chất \(A\).

Vậy bài toán chuyển thành đếm số hoán vị có mọi đoạn hậu tố là tập số nguyên
liên tiếp. Gọi \(g[s,r]\) là số dãy \(C\) tạo bởi các số trong
\(\{s,\ldots,s+r-1\}\), mọi đoạn hậu tố của \(C\) đều là tập số liên tiếp, và
các vị trí đã biết của bài ban đầu vẫn được giữ đúng. Khi đó:
\[
g[s,r]=
\begin{cases}
g[s+1,r-1], & \text{nếu vị trí ngược thứ } r \text{ được xác định là } s,\\
g[s,r-1], & \text{nếu vị trí ngược thứ } r \text{ được xác định là } s+r-1,\\
g[s,r-1]+g[s+1,r-1], & \text{nếu vị trí ngược thứ } r \text{ chưa xác định},\\
0, & \text{trong các trường hợp khác},
\end{cases}
\]
và \(g[s,1]=1\).
Cần tìm \(g[1,n]\).

Nhận xét: điều kiện của bài phức tạp và khó chuyển hóa, rất dễ đẩy thí sinh
vào tìm kiếm. Mấu chốt là quan sát các trường hợp riêng để quy nạp ra hai
phỏng đoán mạnh, chuyển điều kiện phức tạp thành dạng quen thuộc, từ đó mở
đường cho quy hoạch động.

\subsection{Lời kết}

Khi giải bài, ta luôn vô thức phỏng đoán: phỏng đoán tính chất của bài toán,
hình thức của mô hình và nội dung của thuật toán. Phỏng đoán là con đường
quan trọng từ điều đã biết tới điều chưa biết.

Phỏng đoán không phải là ngẫu nhiên cộng may mắn. Nó là suy luận hợp lý nhưng
chủ quan và chưa đầy đủ, xây dựng trên quan sát, phân tích, liên tưởng và quy
nạp. Nó là biểu hiện tổng hợp của năng lực quan sát, khái quát, liên tưởng và
sáng tạo; trình độ phỏng đoán phản ánh chất lượng tư duy của một người.

Theo cách thu được phỏng đoán, có thể chia thành phỏng đoán bằng tương tự và
phỏng đoán bằng quy nạp. Dù dùng phương pháp nào, cũng không thể tách khỏi
quan sát, phân tích và thực hành:
\begin{itemize}
  \item quan sát là dòng máu của phỏng đoán, xuyên suốt toàn bộ quá trình;
  \item phân tích là linh hồn của phỏng đoán, là tiền đề và nền tảng để nêu
  phỏng đoán;
  \item thực hành là bộ khung của phỏng đoán, là nền móng nâng đỡ phỏng đoán;
  rời khỏi thực tế thì phỏng đoán không có ý nghĩa.
\end{itemize}

\end{document}
